\section{Introduction}

Surface electrode arrays placed on neural tissue record local field
potentials (LFPs), reflecting the aggregate synaptic activity of the
underlying neuronal population. When applied to the cerebral cortex
this modality is termed electrocorticography (ECoG); the present system
is developed with implantation on the rat spinal cord as the primary
target application, with the hardware and firmware architecture equally
applicable to other surface recording scenarios.

Compared to scalp EEG, surface electrodes are in much closer proximity
to the target tissue, substantially reducing spatial averaging: the
spatial spread of subdural cortical ECoG can be as local as approximately
0.5--5\,mm~\cite{schwartz2006brain, dubey2019cortical, thielen2023making}.
Recorded amplitudes are correspondingly larger than EEG; for cortical
ECoG, reported values range from hundreds of microvolts to several
millivolts~\cite{thielen2023making} (species unspecified), while
whisker-stimulation-evoked responses in rat cortex have been measured
in the approximate range of 150--850\,\textmu{}V~\cite{buzsaki2012origin,
helmchen1999vivo}. For spinal cord recording, invasive depth electrode
measurements in the rat spinal dorsal horn report C-fibre-evoked field
potentials in the range of approximately 250--500\,\textmu{}V~\cite{sandkuhler1998induction},
and cortically-evoked responses recorded with a penetrating silicon array
in rat spinal cord reach on the order of 300\,\textmu{}V~\cite{song2017spinal};
surface electrode recordings are expected to be lower. The system must
therefore achieve high sensitivity, with the ability to faithfully resolve
signals in the range of tens of microvolts at high signal-to-noise ratio.

LFP signals contain frequency components up to approximately
300\,Hz~\cite{reimann2013biophysically}, which by the Nyquist--Shannon
sampling theorem~\cite{shannon1949} requires a minimum sampling rate
of 600\,Hz per channel. In practice, sampling at the theoretical minimum
introduces reconstruction artefacts and leaves no margin for anti-aliasing
filter roll-off. A sampling rate of 2\,kHz per channel was therefore
chosen, providing more than three times the Nyquist rate and ensuring
high-fidelity waveform reconstruction.

Beyond sufficient sampling rate, temporal regularity of the acquired
samples is critical. Standard spectral analysis methods such as the
discrete Fourier transform (DFT) assume uniformly spaced samples;
irregular sampling intervals introduce spectral artefacts that corrupt
frequency-domain representations. This requirement directly motivates
the hardware-driven acquisition architecture described in
Section~\ref{sec:methods}: by delegating SPI timing entirely to the
PPI subsystem and EasyDMA, sample intervals are determined by a hardware
timer rather than software scheduling, eliminating jitter from CPU
preemption and SoftDevice radio events.
